\section{Introduction}
\label{sec:introduction}

Most tabular recommender portfolios that touch an explicit, multi-level
outcome treat it as an unordered multinomial classification problem, or
collapse it into a binary "good/bad" signal and reach for a ranking metric.
Clothing-fit prediction gives a clean counter-example to that default. A
garment reported as \emph{small}, \emph{fit}, or \emph{large} is not three
unrelated categories: the levels are strictly ordered, and the cost of
confusing \emph{small} for \emph{large} is qualitatively worse than
confusing \emph{small} for \emph{fit}. A model that ignores this ordering
discards information the problem handed it for free.

\subsection{Two contributions, stated precisely}

This report accompanies a working repository, not a novel method. Its two
contributions are methodological:

\begin{enumerate}
  \item \textbf{Ordinal structure, not multinomial.} We fit a
        proportional-odds cumulative link model (Section~\ref{sec:methodology})
        instead of treating the three fit levels as unordered classes, and we
        show this pays off on ordinal-aware metrics
        (Section~\ref{sec:results}) even though both models see identical
        features.
  \item \textbf{Calibrated uncertainty, not a bare label.} Wrapping the
        fitted model in split conformal prediction turns a point guess into
        a \emph{set} of plausible outcomes with a finite-sample marginal
        coverage guarantee (Section~\ref{sec:methodology}). For a sizing
        recommendation, "this could run fit or large" is a more honest and
        more actionable signal than a single confident label that is
        sometimes simply wrong.
\end{enumerate}

Both pieces are classical machine learning: the cumulative link model is fit
by maximizing a smooth, explicit likelihood with
\texttt{scipy.optimize.minimize}, and the conformal calibration is closed-form
arithmetic over a held-out calibration set. No neural network, embedding
learned by gradient descent, or deep framework appears anywhere in the
pipeline.

\subsection{Evaluation honesty}

Users are split at the group level: every review by a given user lands in
exactly one of train, validation, or test
(Section~\ref{sec:data}). This prevents the model from memorizing one user's
exact body-measurement combination and matching it back to that user's
historical fit outcomes -- a leakage mode specific to supervised tabular
problems with repeated subjects, distinct from the temporal point-in-time
leakage that recommender systems usually worry about. Both leakage modes are
present in this dataset and are tested independently
(Section~\ref{sec:feature-engineering}).

Hyperparameters for every model, baselines included, are selected on the
validation partition; the test partition is read once, to produce
Table~\ref{tab:results}.

\subsection{Repository}

The implementation and the scripts that generate every table in this
document are available in the accompanying repository. The results table
(Table~\ref{tab:results}) is emitted by
\texttt{scripts/export\_results\_table.py} and the feature catalogue
(Table~\ref{tab:features}) by \texttt{scripts/export\_feature\_tables.py}, so
the document cannot drift away from the code.
